\documentclass[sigconf,anonymous,review,natbib=false]{acmart}

\AtBeginDocument{\providecommand\BibTeX{{Bib\TeX}}}
\setcopyright{none}
\settopmatter{printacmref=false,printfolios=true}
\renewcommand\footnotetextcopyrightpermission[1]{}
\acmYear{2027}
\acmConference[SAC '27]{The 42nd ACM/SIGAPP Symposium on Applied Computing}{April 5--9, 2027}{Gwangju, South Korea}

\RequirePackage[datamodel=acmdatamodel,style=acmnumeric]{biblatex}
\addbibresource{movie_recommender_references.bib}

\begin{document}

\title{When Incremental Feedback Hurts: A Reproducible Synthetic Evaluation of Cold-Start Movie Recommendation}
\renewcommand{\shorttitle}{When Incremental Feedback Hurts}

\begin{abstract}
Conversational recommender systems can elicit initial preferences and then
adapt their rankings from explicit feedback. However, adaptation does not
necessarily improve recommendation quality, especially when feedback is scarce
and the initial profile is already informative. This paper presents a
reproducible evaluation of incremental personalization for cold-start movie
recommendation. We compare six methods, from popularity and content-based
baselines to static-profile and incremental recommenders, across 15,120 frozen
experimental cells generated from 35 independent synthetic agents. The primary
paired analysis compares the incremental method with its static-profile
control using NDCG@5. Incremental personalization decreased NDCG@5 by 0.0181
(3.25\% relative to the static baseline; 95\% bootstrap CI
$[-0.0343,-0.0022]$; two-sided Wilcoxon $p=0.0356$; rank-biserial
correlation $=-0.406$). Twenty-six agents lost quality and nine improved.
Performance did not increase monotonically with additional feedback, and a
Friedman test across feedback conditions was not significant ($p=0.0960$).
Exploratory subgroup results indicate larger losses for richer initial
profiles. These findings do not establish effects on human users; they show
that, under a controlled and versioned simulator, an unregularized feedback
update can displace a strong initial profile and harm ranking quality. A preregistered external evaluation on MovieLens 100K retained 43 eligible
users: B5 and B4 tied for every user, while B5 underperformed popularity by
0.0785 NDCG@5. This partial, semantically adapted validation neither replicates
nor reverses the synthetic effect. The work contributes traceable infrastructure,
negative evidence, and an explicit boundary on transportability.
\end{abstract}

\begin{CCSXML}
<ccs2012>
 <concept>
  <concept_id>10002951.10003260.10003277</concept_id>
  <concept_desc>Information systems~Recommender systems</concept_desc>
  <concept_significance>500</concept_significance>
 </concept>
 <concept>
  <concept_id>10002951.10003260.10003309</concept_id>
  <concept_desc>Information systems~Evaluation of retrieval results</concept_desc>
  <concept_significance>300</concept_significance>
 </concept>
 <concept>
  <concept_id>10010147.10010257</concept_id>
  <concept_desc>Computing methodologies~Machine learning</concept_desc>
  <concept_significance>100</concept_significance>
 </concept>
</ccs2012>
\end{CCSXML}

\ccsdesc[500]{Information systems~Recommender systems}
\ccsdesc[300]{Information systems~Evaluation of retrieval results}
\ccsdesc[100]{Computing methodologies~Machine learning}

\keywords{recommender systems, cold start, incremental personalization,
synthetic agents, reproducibility}

\maketitle

\section{Introduction}

Recommender systems reduce information overload by ranking items according to
estimated user preferences~\cite{ricci2015recommender}. A persistent challenge
is user cold start: before interaction history exists, collaborative evidence
is unavailable and personalization must rely on limited elicited information
or non-personalized priors~\cite{schein2002methods}. Conversational interfaces
offer a practical response by collecting preferences and explicit feedback
during interaction~\cite{jannach2021survey}. Yet the existence of feedback does
not imply that every update policy improves the next ranking. Sparse,
contradictory, or exposure-biased observations may move the representation away
from an already useful declared profile.

This paper studies that failure mode in a movie recommender. A user starts with
up to three ranked genres, a decade preference, and a popularity preference.
The incremental method updates this representation after positive and negative
feedback. We ask whether this update improves ranking quality over the same
recommender with a static profile, how quality changes with more feedback, and
how the effect varies across cold-start profiles and simulated behaviors.

The study makes three contributions. First, it provides a common ranking
contract and a frozen, hash-validated matrix covering six methods, six feedback
conditions, six cold-start profiles, seven synthetic personas, five seeds, and
two list sizes. Second, it performs a paired analysis at the independent-agent
level rather than treating repeated rankings as independent observations.
Third, it preserves a negative result: incremental feedback significantly
reduced NDCG@5 relative to a static-profile control in the simulated setting.
The claim is deliberately bounded to the generator and catalog used; no human
preference or satisfaction claim is made.

\section{Background and Related Work}

Content-based recommendation represents users and items in a shared feature
space, making it suitable when interaction histories are limited
~\cite{lops2011content}. Hybrid systems combine complementary signals to reduce
the limitations of individual recommenders~\cite{burke2002hybrid}. Evaluation,
however, must match the task: classification accuracy over constructed labels
does not replace ranking evaluation, and offline metrics do not directly imply
user utility~\cite{herlocker2004evaluating}.

Cold-start studies have long emphasized the distinction between accuracy and
the ability to learn useful preferences from little evidence
~\cite{schein2002methods}. Conversational recommendation adds preference
elicitation and interaction, but also introduces choices about question policy,
feedback interpretation, and online updates~\cite{jannach2021survey}. Beyond
relevance, recommender evaluation may consider diversity, novelty, coverage,
and popularity effects~\cite{castells2022novelty}. Our analyses report ranking, discovery, robustness, ablation, and cost outcomes,
while preserving their confirmatory or exploratory status.

\section{Research Questions}

The primary question (RQ1) asks whether incremental positive and negative
feedback improves NDCG@5 over a static initial profile. H1 predicts a positive
paired difference between the incremental method B5 and static method B4. RQ2
asks whether ranking quality increases as the available history grows from zero
to one, three, five, ten, and all prior interactions (H2). Exploratory questions
consider cold-start information, popularity, discovery, and component
contribution. Only RQ1/H1 is treated as confirmatory in this version.

\section{Methodology}

\subsection{Catalog, profiles, and agents}

The experiment uses a frozen movie catalog with structured and textual
metadata. Candidate identity and eligibility are shared by all methods in a
comparison. Cold-start profiles P0--P5 progressively expose no declared
preference, one genre, three ranked genres, decade, popularity, and available
feedback. Synthetic agents implement seven versioned behavioral personas,
including consistent, noisy, contradictory, exploratory, broad, specific, and
popularity-sensitive behaviors. Each persona is instantiated under five seeds,
yielding 35 independent longitudinal units.

Synthetic agents are used because the available real feedback is insufficient
for inference. Their latent preferences generate controlled relevance and
feedback. This enables reproducible stress testing but creates construct and
external-validity limitations: agents are not people, personas are not
demographic groups, and shared feature assumptions can introduce circularity.

\subsection{Compared methods}

We evaluate B0 popularity, B1 content-based similarity, B2 TF--IDF cosine
similarity, B3 supervised relevance, B4 static profile, and B5 incremental
personalization. B4 and B5 share the same base representation and candidates;
B5 alone incorporates feedback available before the current ranking. A planned
hybrid B6 was excluded before holdout evaluation because no distinct,
validation-selected mechanism had been frozen.

\subsection{Experimental design and reproducibility}

The full factorial matrix contains 15,120 cells: six methods, six feedback
conditions C0--C5, six profiles P0--P5, seven personas, five seeds, and
$K\in\{5,10\}$. Every cell records method and protocol versions, seed, candidate
identity, logical time, metrics, and status. Manifests and cell hashes prevent
mixing pilot, interrupted, or official executions. Code, configuration, DVC
pointers, and derived-report commands are versioned; credentials and sensitive
or unauthorized data are excluded.

\subsection{Metrics and statistical analysis}

NDCG@5 is the primary metric. For H1, observations from C1--C5 and P0--P5 are
first averaged within each agent and method. C0 is excluded because B5 has not
received feedback and is contractually equivalent to B4. The 35 resulting B5
and B4 values are compared by their paired difference. We report the mean,
median, a paired bootstrap 95\% confidence interval with 10,000 resamples and
seed 42, a two-sided Wilcoxon signed-rank test, and paired rank-biserial
correlation. H1 requires a positive difference, compatible confidence interval,
the pre-defined significance threshold, and a reported effect size.

For RQ2, B5 is compared across C0--C5 within agents by a Friedman test.
Post-hoc comparisons against C0 are allowed only if the global test rejects the
null hypothesis. Subgroup analyses by condition, profile, persona, seed, and K
are exploratory and use Holm correction within each dimension and K.

\section{Results}

\subsection{Primary comparison}

Table~\ref{tab:primary} reports the confirmatory result. B5 obtained lower mean
NDCG@5 than B4. The mean paired difference was $-0.0181$, or $-3.25\%$ relative
to B4. Its confidence interval excluded zero, and the effect was negative and
moderate by rank-biserial correlation. H1 is therefore not supported; the data
instead provide evidence of harm within the synthetic environment.

\begin{table}[t]
  \caption{Primary paired analysis at the synthetic-agent level.}
  \label{tab:primary}
  \begin{tabular}{lr}
    \toprule
    Quantity & Value \\
    \midrule
    Independent agents & 35 \\
    Mean NDCG@5, B4 & 0.5562 \\
    Mean NDCG@5, B5 & 0.5381 \\
    Mean B5 $-$ B4 & $-0.0181$ \\
    Bootstrap 95\% CI & $[-0.0343,-0.0022]$ \\
    Wilcoxon $p$ (two-sided) & 0.0356 \\
    Rank-biserial correlation & $-0.4063$ \\
    Gains / losses / ties & 9 / 26 / 0 \\
    \bottomrule
  \end{tabular}
\end{table}

\subsection{Feedback quantity and cold start}

B5 mean NDCG@5 relative to C0 changed by $-0.0131$, $-0.0179$, $-0.0161$,
$-0.0158$, and $+0.0107$ at C1--C5, respectively. The trajectory is not
monotonic. The Friedman test did not reject equal performance across conditions
($\chi^2=9.3473$, $p=0.0960$), so post-hoc tests were not performed and H2 is
not supported.

Cold-start summaries show a substantial descriptive difference between P0 and
P5 (B5: 0.3251 versus 0.6271; B4: 0.2926 versus 0.6738). A pre-specified paired contrast found P5--P0 under B5 equal to 0.2819
(Holm-adjusted $p=2.33\times10^{-10}$; rank-biserial $=0.9968$),
supporting H3 only inside the simulator.

\subsection{Exploratory heterogeneity}

After Holm correction, negative B5--B4 effects remained for C3 and profiles
P3--P5 at K=5, and P3 at K=10. For P5/K=5, the difference was $-0.0560$
(95\% CI $[-0.0871,-0.0267]$; raw $p=0.00183$; Holm $p=0.00930$;
rank-biserial $=-0.6411$). Only seed 42 favored B5 on average; the other four
seeds were negative. Persona and seed estimates have only five and seven units,
respectively, and are treated as imprecise diagnostics.

Across methods, descriptive mean NDCG@5 was 0.5549 for B1/B4, 0.5399 for B5,
0.5230 for B3, 0.4598 for B2, and 0.3471 for B0. B5 exceeded B0 by 0.1898 NDCG@5 in the synthetic paired contrast and increased
novelty and coverage, at a higher median latency (80.19 versus 18.70 ms).

\section{Discussion}

The result challenges the assumption that more interaction automatically means
better recommendation. B4 already receives a structured profile. An
unregularized update from sparse feedback can overreact to a small exposed set
and displace useful prior information. Larger exploratory losses for P3--P5 are
consistent with this mechanism: the richer the initial profile, the more there
is to damage. Validation-only ablations support this interpretation: removing the complete
history increased NDCG@5 by 0.0380, whereas removing genre reduced it by 0.1656.
Contradictory-feedback stress tests produced a 0.2077 loss relative to B4; these
remain simulator-specific mechanism checks, not human causal effects.

The statistically detectable average loss is practically modest, but its
direction is consistent for 26 of 35 agents and its effect size is not trivial.
The correct engineering response is not to discard feedback; it is to gate or
regularize updates, model confidence, and test contradictory and random signals.
The negative result also illustrates why near-perfect classification against a
constructed relevance proxy would be insufficient: the downstream ranking can
still deteriorate after adaptation.

\subsection{External validation}

We evaluated transportability on MovieLens 100K~\cite{harper2015movielens}
using a frozen global chronological 70/10/20 split. Of 943 users, 43 met the
pre-specified history and relevance criteria. B5--B4 was exactly 0.0000
(95\% CI [0.0000, 0.0000]; 43 ties), whereas B5--B0 was -0.0785
(95\% CI [-0.1289, -0.0341]; $p=0.00428$). The external B4/B5 variants
use genre histories rather than declared profiles and simulated conditions, so
the neutral result is classified as inconclusive partial transportability.
Synthetic and public samples were never pooled.

\section{Threats to Validity}

\textbf{Internal validity.} Pairing and hash validation reduce execution mixing
and pseudo-replication. Validation-only ablations and robustness runs isolate mechanisms without
reusing the official holdout, but they do not establish human causality.

\textbf{Construct validity.} Synthetic NDCG measures alignment with latent
simulated preferences, not satisfaction or perceived usefulness. Method,
feedback, and relevance share a feature domain, creating potential circularity.
Novelty and popularity bias are unavailable, and no diversity non-inferiority
margin was frozen.

\textbf{External validity.} The filtered catalog and seven personas do not
represent the full variety of people, contexts, and preference drift. Real feedback is scarce and no prospective human experiment was performed.
MovieLens adds observed ratings but only partial construct compatibility and low
eligible-user coverage. Findings cannot be generalized to Telegram users.

\textbf{Conclusion validity.} The primary analysis respects the independent
unit and reports uncertainty and effect size. Subgroups remain exploratory even
after multiplicity correction. Missing studies are reported as unavailable,
not as null effects.

\section{Conclusion and Next Steps}

In a reproducible synthetic experiment, incremental personalization reduced
NDCG@5 relative to a static-profile recommender and did not improve monotonically
with more feedback. The contribution is therefore not a claim that B5 is a
better recommender. It is an auditable evaluation showing that naive adaptation
can harm a strong cold-start profile and that repeated observations must be
analyzed at the agent level.

Ablation, robustness, confirmatory contrasts, discovery and cost analyses, and
a partial external validation are now complete. The evidence supports the
narrow conclusion that naive adaptation can be harmful or inert across the two
evaluated settings; it does not establish human utility. Future work should
validate a regularized mechanism on denser public histories and, when feasible,
under an authorized human protocol.

\printbibliography

\end{document}
